\documentclass[11pt]{article}
\usepackage{amsmath, amssymb, amsthm}
\usepackage{array}
\usepackage{booktabs}
\usepackage[margin=0.75in]{geometry}
\usepackage{xcolor}
\usepackage{tcolorbox}
\tcbuselibrary{breakable,theorems,skins}
\usepackage{enumitem}
\usepackage{fancyhdr}
\usepackage{graphicx}
\usepackage{multicol}
\usepackage{tikz}
\usetikzlibrary{shapes, arrows, positioning, calc}
\usepackage[utf8]{inputenc}
\usepackage[
    colorlinks=true,
    linkcolor=blue,
    urlcolor=blue,
    citecolor=blue,
    bookmarks=true,
    bookmarksnumbered=true,
    pdfpagemode=UseOutlines,
    pdftitle={Statistical Foundation - Exam Formula Sheet},
    pdfauthor={CSU1658},
    pdfsubject={Statistical Foundation of Data Science},
    pdfkeywords={statistics, formulas, data science, exam}
]{hyperref}

\definecolor{formulabox}{RGB}{230,240,255}
\definecolor{importantbox}{RGB}{255,240,240}
\definecolor{examplebox}{RGB}{240,255,240}

\newtcolorbox{formula}[1]{
    colback=formulabox,
    colframe=blue!75!black,
    title=#1,
    fonttitle=\bfseries,
    breakable
}

\newtcolorbox{important}[1]{
    colback=importantbox,
    colframe=red!75!black,
    title=\textbf{#1} (HIGH PRIORITY),
    fonttitle=\bfseries,
    breakable
}

\newtcolorbox{example}{
    colback=examplebox,
    colframe=green!75!black,
    title=Quick Reference,
    fonttitle=\bfseries,
    breakable
}

\pagestyle{fancy}
\fancyhf{}
\lhead{Statistical Foundation - Exam Formulas}
\rhead{CSU1658}
\cfoot{\thepage}
\setlength{\headheight}{14pt}

\setlength{\parindent}{0pt}
\setlength{\parskip}{6pt}

\title{\textbf{EXAM FORMULA SHEET}\\
\large Statistical Foundation of Data Science\\
Quick Reference Guide}
\author{CSU1658}
\date{December 2025}

\begin{document}

\maketitle
\tableofcontents
\newpage

% ============================================================
% DESCRIPTIVE STATISTICS
% ============================================================

\section{Descriptive Statistics}

\begin{important}{Mean (Average) - HIGH PROBABILITY}
\textbf{Population Mean:}
\[\mu = \frac{1}{N}\sum_{i=1}^{N} x_{i} = \frac{x_1 + x_2 + \cdots + x_N}{N}\]

\textbf{Sample Mean:}
\[\bar{x} = \frac{1}{n}\sum_{i=1}^{n} x_{i} = \frac{x_1 + x_2 + \cdots + x_n}{n}\]

\textbf{Weighted Mean:}
\[\bar{x}_w = \frac{\sum_{i=1}^{n} w_i x_{i}}{\sum_{i=1}^{n} w_i} = \frac{w_1x_1 + w_2x_2 + \cdots + w_nx_n}{w_1 + w_2 + \cdots + w_n}\]

\textbf{Geometric Mean:}
\[\bar{x}_{\text{geo}} = \sqrt[n]{x_1 \cdot x_2 \cdot \ldots \cdot x_n} = \left(\prod_{i=1}^{n} x_i\right)^{1/n}\]

\textbf{Harmonic Mean:}
\[\bar{x}_{\text{harm}} = \frac{n}{\sum_{i=1}^{n} \frac{1}{x_i}} = \frac{n}{\frac{1}{x_1} + \frac{1}{x_2} + \cdots + \frac{1}{x_n}}\]

\textbf{Properties:} $\mathbb{E}[\bar{X}] = \mu$, $\text{Var}(\bar{X}) = \frac{\sigma^2}{n}$, $\text{SE}(\bar{X}) = \frac{\sigma}{\sqrt{n}}$

\textbf{When to use:} Arithmetic mean for symmetric data; Geometric for growth rates; Harmonic for rates/ratios.
\end{important}

\begin{formula}{Median - MEDIUM PROBABILITY}
\textbf{Ordered dataset:} Sort data in ascending order

\textbf{For odd n:}
\[\text{Median} = x_{\left(\frac{n+1}{2}\right)}\]

\textbf{For even n:}
\[\text{Median} = \frac{x_{\left(\frac{n}{2}\right)} + x_{\left(\frac{n}{2}+1\right)}}{2}\]

\textbf{When to use:} Better than mean for skewed distributions or data with outliers. Not affected by extreme values.
\end{formula}

\begin{formula}{Mode - LOW PROBABILITY}
\textbf{Definition:} The value that appears most frequently in the dataset.

\textbf{Properties:}
- Unimodal: One mode
- Bimodal: Two modes
- Multimodal: Multiple modes
- No mode: All values appear with same frequency

\textbf{When to use:} Categorical data, finding most common category.
\end{formula}

\begin{important}{Variance and Standard Deviation - VERY HIGH PROBABILITY}
\textbf{Population Variance:}
\[\sigma^2 = \frac{1}{N}\sum_{i=1}^{N}(x_{i} - \mu)^2 = \mathbb{E}[(X - \mu)^2] = \mathbb{E}[X^2] - (\mathbb{E}[X])^2\]

\textbf{Sample Variance (Unbiased Estimator):}
\[s^2 = \frac{1}{n-1}\sum_{i=1}^{n}(x_{i} - \bar{x})^2 = \frac{\sum(x_i - \bar{x})^2}{n-1}\]

\textbf{Computational Formula (Shortcut):}
\[s^2 = \frac{1}{n-1}\left(\sum_{i=1}^{n}x_i^2 - \frac{\left(\sum_{i=1}^{n}x_i\right)^2}{n}\right) = \frac{n\sum x_i^2 - \left(\sum x_i\right)^2}{n(n-1)}\]

\textbf{Population Standard Deviation:}
\[\sigma = \sqrt{\sigma^2}, \quad \text{SD}[X] = \sqrt{\text{Var}[X]}\]

\textbf{Sample Standard Deviation:}
\[s = \sqrt{s^2}, \quad \text{SE}(\bar{x}) = \frac{s}{\sqrt{n}}\]

\textbf{Coefficient of Variation (Relative Dispersion):}
\[\text{CV} = \frac{s}{\bar{x}} \times 100\% = \frac{\sigma}{\mu} \times 100\%\]

\textbf{Properties of Variance:}
\begin{itemize}[nosep]
\item $\text{Var}(aX + b) = a^2\text{Var}(X)$ (linear transformation)
\item $\text{Var}(X + Y) = \text{Var}(X) + \text{Var}(Y) + 2\text{Cov}(X,Y)$
\item $\text{Var}(X - Y) = \text{Var}(X) + \text{Var}(Y) - 2\text{Cov}(X,Y)$
\item If $X \perp Y$ (independent): $\text{Var}(X \pm Y) = \text{Var}(X) + \text{Var}(Y)$
\item $\text{Var}\left(\sum_{i=1}^{n}X_i\right) = \sum_{i=1}^{n}\text{Var}(X_i) + 2\sum_{i<j}\text{Cov}(X_i,X_j)$
\end{itemize}

\textbf{Bessel's Correction:} Division by $(n-1)$ instead of $n$ makes $s^2$ unbiased: $\mathbb{E}[s^2] = \sigma^2$.

\textbf{Interpretation:} Measures spread/dispersion around mean. Larger $\sigma^2$ = more variability.
\end{important}

\begin{important}{Standard Error (SE) - HIGH PROBABILITY}
\textbf{Standard Error of Mean:}
\[\text{SE}(\bar{x}) = \frac{\sigma}{\sqrt{n}} \approx \frac{s}{\sqrt{n}}\]

\textbf{Standard Error of Proportion:}
\[\text{SE}(\hat{p}) = \sqrt{\frac{p(1-p)}{n}} \approx \sqrt{\frac{\hat{p}(1-\hat{p})}{n}}\]

\textbf{Standard Error of Difference of Means:}
\[\text{SE}(\bar{x}_1 - \bar{x}_2) = \sqrt{\frac{\sigma_1^2}{n_1} + \frac{\sigma_2^2}{n_2}} \approx \sqrt{\frac{s_1^2}{n_1} + \frac{s_2^2}{n_2}}\]

\textbf{Interpretation:} Measures precision of sample statistic as estimator of population parameter. Smaller SE $\Rightarrow$ more precise estimate.
\end{important}

\begin{important}{Standard Deviation - HIGH PROBABILITY}
\textbf{Population SD:}
\[\sigma = \sqrt{\sigma^2} = \sqrt{\frac{1}{N}\sum_{i=1}^{N}(x_{i} - \mu)^2}\]

\textbf{Sample SD:}
\[s = \sqrt{s^2} = \sqrt{\frac{1}{n-1}\sum_{i=1}^{n}(x_{i} - \bar{x})^2}\]

\textbf{Coefficient of Variation (relative SD):}
\[CV = \frac{s}{\bar{x}} \times 100\%\]

\textbf{Interpretation:} Average distance from mean (same units as data). Smaller SD = less spread.
\end{important}
\textbf{Frequency:} Count of occurrences of each value
\[f_i = \text{count of } x_{i}\]

\textbf{Relative Frequency:}
\[RF_i = \frac{f_i}{n}\]

\textbf{Cumulative Frequency:}
\[CF_i = \sum_{j=1}^{i} f_j\]

\textbf{Cumulative Relative Frequency:}
\[CRF_i = \frac{CF_i}{n}\]
\end{formula}

\begin{formula}{Range IQR and Percentiles}
\textbf{Range:}
\[\text{Range} = x_{\max} - x_{\min}\]

\textbf{Percentiles:} Value below which P\% of data falls
\[P_k = \text{value at position } \frac{k(n+1)}{100}\]

\textbf{Quartiles:}
- Q$_1$ (25th percentile): 1st quartile
- Q$_2$ (50th percentile): Median
- Q$_3$ (75th percentile): 3rd quartile

\textbf{Interquartile Range (IQR):}
\[IQR = Q_3 - Q_1\]

\textbf{Outlier Detection with IQR:}
- Lower fence: Q$_1$ - 1.5$\times$IQR
- Upper fence: Q$_3$ + 1.5$\times$IQR
- Outliers: Values outside fences
\end{formula}

\begin{formula}{Skewness and Kurtosis}
\textbf{Skewness (Fisher's moment coefficient):}
\[g_1 = \frac{n}{(n-1)(n-2)}\sum_{i=1}^{n}\left(\frac{x_{i} - \bar{x}}{s}\right)^3\]

\textbf{Interpretation:}
- g$_1$ = 0: Symmetric distribution
- g$_1$ > 0: Right-skewed (positive skew)
- g$_1$ < 0: Left-skewed (negative skew)

\textbf{Kurtosis (excess kurtosis):}
\[g_2 = \left[\frac{n(n+1)}{(n-1)(n-2)(n-3)}\sum_{i=1}^{n}\left(\frac{x_{i} - \bar{x}}{s}\right)^4\right] - \frac{3(n-1)^2}{(n-2)(n-3)}\]

\textbf{Interpretation:}
- g$_2$ = 0: Mesokurtic (normal)
- g$_2$ > 0: Leptokurtic (heavy tails)
- g$_2$ < 0: Platykurtic (light tails)
\end{formula}

\begin{important}{Covariance and Correlation - VERY HIGH PROBABILITY}
\textbf{Covariance:}
\[\text{Cov}(X,Y) = \frac{1}{n-1}\sum_{i=1}^{n}(x_{i} - \bar{x})(y_{i} - \bar{y})\]

\textbf{Pearson Correlation Coefficient:}
\[r = \frac{\text{Cov}(X,Y)}{s_X s_Y} = \frac{\sum(x_{i} - \bar{x})(y_{i} - \bar{y})}{\sqrt{\sum(x_{i} - \bar{x})^2 \sum(y_{i} - \bar{y})^2}}\]

\textbf{Properties:}
- Range: -1 $\leq$ r $\leq$ 1
- r = 1: Perfect positive correlation
- r = -1: Perfect negative correlation
- r = 0: No linear correlation

\textbf{Coefficient of Determination:}
\[R^2 = r^2\]
Proportion of variance in Y explained by X.

\textbf{Spearman Rank Correlation:}
\[r_s = 1 - \frac{6\sum $d_{i}$^2}{n(n^2-1)}\]
where $d_{i}$ = rank difference for observation i.

\textbf{When to use:}
- Pearson: Linear relationships, normally distributed
- Spearman: Monotonic relationships, ordinal data, non-normal
\end{important}

\begin{example}
\textbf{Quick Tips:}
- Mean is sensitive to outliers; Median is robust
- SD measures spread; larger SD = more variability
- Always use n-1 for sample variance (Bessel's correction)
- Variance units are squared; SD has same units as data
- IQR detects outliers: beyond Q$_1$-1.5$\times$IQR or Q$_3$+1.5$\times$IQR
- Correlation measures linear association, not causation
\end{example}

% ============================================================
% PROBABILITY DISTRIBUTIONS
% ============================================================

\section{Probability Distributions}

\subsection{PMF and PDF}

\begin{important}{Probability Mass Function (PMF) - Discrete - HIGH PROBABILITY}
\textbf{Definition:} For discrete random variable X
\[P(X = x) = p(x)\]

\textbf{Properties:}
\begin{align*}
0 &\leq p(x) \leq 1 \text{ for all } x \\
\sum_{\text{all } x} p(x) &= 1
\end{align*}

\textbf{Expected Value:}
\[\mathbb{E}[X] = \sum_{\text{all } x} x \cdot p(x)\]

\textbf{Variance:}
\[\text{Var}(X) = \mathbb{E}[X^2] - (\mathbb{E}[X])^2 = \sum x^2 p(x) - \mu^2\]
\end{important}

\begin{important}{Probability Density Function (PDF) - Continuous - HIGH PROBABILITY}
\textbf{Definition:} For continuous random variable X
\[f(x) \geq 0 \text{ for all } x\]

\textbf{Normalization:}
\[\int_{-\infty}^{\infty} f(x)\,dx = 1\]

\textbf{Probability:}
\[P(a \leq X \leq b) = \int_a^b f(x)\,dx\]

\textbf{Expected Value:}
\[\mathbb{E}[X] = \int_{-\infty}^{\infty} x \cdot f(x)\,dx\]

\textbf{Variance:}
\[\text{Var}(X) = \int_{-\infty}^{\infty} (x-\mu)^2 f(x)\,dx\]
\end{important}

\subsection{Bernoulli Distribution}

\begin{formula}{Bernoulli - MEDIUM PROBABILITY}
\textbf{PMF:} Single trial with success probability p
\[P(X = x) = \begin{cases}
p & \text{if } x = 1 \\
1-p & \text{if } x = 0
\end{cases}\]

Or: \(P(X=x) = p^x(1-p)^{1-x}\) for x $\in$ {0,1}

\textbf{Expected Value:}
\[\mathbb{E}[X] = p\]

\textbf{Variance:}
\[\text{Var}(X) = p(1-p)\]

\textbf{When to use:} Single yes/no trial, coin flip, binary outcome.
\end{formula}

\subsection{Binomial Distribution}

\begin{important}{Binomial Distribution - HIGH PROBABILITY}
\textbf{PMF:} n independent trials, probability p of success
\[P(X = k) = \binom{n}{k}p^k(1-p)^{n-k}\]

where \(\binom{n}{k} = \frac{n!}{k!(n-k)!}\)

\textbf{Expected Value:}
\[\mathbb{E}[X] = np\]

\textbf{Variance:}
\[\text{Var}(X) = np(1-p)\]

\textbf{Standard Deviation:}
\[\sigma = \sqrt{np(1-p)}\]

\textbf{When to use:} Fixed number of independent trials, each with same success probability. Count of successes.
\end{important}

\subsection{Poisson Distribution}

\begin{important}{Poisson Distribution - HIGH PROBABILITY}
\textbf{PMF:} Events occurring at rate $\lambda$ per unit time
\[P(X = k) = \frac{e^{-\lambda}\lambda^k}{k!}, \quad k = 0,1,2,\ldots\]

\textbf{Expected Value:}
\[\mathbb{E}[X] = \lambda\]

\textbf{Variance:}
\[\text{Var}(X) = \lambda\]

\textbf{Standard Deviation:}
\[\sigma = \sqrt{\lambda}\]

\textbf{When to use:} Rare events, arrivals per unit time, events in fixed interval.

\textbf{Poisson Approximation to Binomial:} When n large, p small, np = $\lambda$ moderate.
\end{important}

\begin{formula}{Central Limit Theorem (CLT) - MEDIUM PROBABILITY}
\textbf{Statement:} For large n, the sampling distribution of the sample mean approaches normal distribution.

\[\bar{X} \sim N\left(\mu, \frac{\sigma^2}{n}\right)\]

\textbf{Standardized form:}
\[Z = \frac{\bar{X} - \mu}{\sigma/\sqrt{n}} \sim N(0,1)\]

\textbf{Requirements:}
- Sample size n $\geq$ 30 (rule of thumb)
- Independent observations
- Identically distributed

\textbf{Applications:} Hypothesis testing, confidence intervals, approximations
\end{formula}

\begin{formula}{Law of Large Numbers (LLN)}
\textbf{Weak LLN:} As n → $\infty$, sample mean converges in probability to population mean.

\[\lim_{n \to \infty} P(|\bar{X}_n - \mu| > \epsilon) = 0\]

\textbf{Practical meaning:} Larger samples give better estimates.
\end{formula}

\begin{important}{Normal Distribution $\mathcal{N}(\mu, \sigma^2)$ - VERY HIGH PROBABILITY}
\textbf{Probability Density Function:}
\[f(x) = \frac{1}{\sigma\sqrt{2\pi}} \exp\left(-\frac{(x-\mu)^2}{2\sigma^2}\right) = \frac{1}{\sigma\sqrt{2\pi}} e^{-\frac{1}{2}\left(\frac{x-\mu}{\sigma}\right)^2}, \quad x \in \mathbb{R}\]

\textbf{Standard Normal $\mathcal{N}(0,1)$:}
\[\phi(z) = \frac{1}{\sqrt{2\pi}} e^{-z^2/2}, \quad z \in \mathbb{R}\]

\textbf{Cumulative Distribution Function:}
\[\Phi(z) = P(Z \leq z) = \int_{-\infty}^{z} \phi(t)\,dt = \int_{-\infty}^{z} \frac{1}{\sqrt{2\pi}} e^{-t^2/2}\,dt\]

\textbf{Standardization (Z-transformation):}
\[Z = \frac{X - \mu}{\sigma} \sim \mathcal{N}(0,1), \quad P(a \leq X \leq b) = \Phi\left(\frac{b-\mu}{\sigma}\right) - \Phi\left(\frac{a-\mu}{\sigma}\right)\]

\textbf{Parameters:} $\mathbb{E}[X] = \mu$, $\text{Var}(X) = \sigma^2$, Skewness $= 0$, Excess Kurtosis $= 0$

\textbf{Empirical Rule (68-95-99.7 Rule):}
\begin{itemize}[nosep]
\item $P(\mu - \sigma \leq X \leq \mu + \sigma) \approx 0.6827$ (68.27\%)
\item $P(\mu - 2\sigma \leq X \leq \mu + 2\sigma) \approx 0.9545$ (95.45\%)
\item $P(\mu - 3\sigma \leq X \leq \mu + 3\sigma) \approx 0.9973$ (99.73\%)
\end{itemize}

\textbf{Properties:}
\begin{itemize}[nosep]
\item Symmetric about $\mu$: $f(\mu + x) = f(\mu - x)$
\item Bell-shaped, unimodal
\item $\Phi(-z) = 1 - \Phi(z)$ (symmetry of CDF)
\item Linear transformation: If $X \sim \mathcal{N}(\mu, \sigma^2)$, then $aX + b \sim \mathcal{N}(a\mu + b, a^2\sigma^2)$
\item Sum of independent normals: $X_1 \sim \mathcal{N}(\mu_1, \sigma_1^2), X_2 \sim \mathcal{N}(\mu_2, \sigma_2^2)$ \\
$\Rightarrow X_1 + X_2 \sim \mathcal{N}(\mu_1 + \mu_2, \sigma_1^2 + \sigma_2^2)$
\end{itemize}

\textbf{Critical Values (Common):}
- $z_{0.10} = 1.28$, $z_{0.05} = 1.645$, $z_{0.025} = 1.96$, $z_{0.01} = 2.33$, $z_{0.005} = 2.576$
\end{important}

\subsection{Exponential Distribution}

\begin{formula}{Exponential Distribution - LOW PROBABILITY}
\textbf{PDF:} For rate parameter $\lambda$ > 0
\[f(x) = \lambda e^{-\lambda x}, \quad x \geq 0\]

\textbf{CDF:}
\[F(x) = 1 - e^{-\lambda x}\]

\textbf{Expected Value:}
\[\mathbb{E}[X] = \frac{1}{\lambda}\]

\textbf{Variance:}
\[\text{Var}(X) = \frac{1}{\lambda^2}\]

\textbf{Memoryless Property:}
\[P(X > s + t | X > s) = P(X > t)\]

\textbf{When to use:} Time between events, waiting times, lifetimes
\end{formula}

\begin{important}{Normal Distribution - VERY HIGH PROBABILITY}
\textbf{PDF:} X $\sim$ N($\mu$, $\sigma$²)
\[f(x) = \frac{1}{\sigma\sqrt{2\pi}}\exp\left(-\frac{(x-\mu)^2}{2\sigma^2}\right)\]

\textbf{Standard Normal:} Z $\sim$ N(0,1)
\[f(z) = \frac{1}{\sqrt{2\pi}}e^{-z^2/2}\]

\textbf{Standardization (Z-score):}
\[Z = \frac{X - \mu}{\sigma}\]

\textbf{Properties:}
- Mean = Median = Mode = $\mu$
- 68\% within $\mu$ ± $\sigma$
- 95\% within $\mu$ ± 1.96$\sigma$
- 99.7\% within $\mu$ ± 3$\sigma$ (Empirical Rule)

\textbf{Linear Transformation:} If X $\sim$ N($\mu$, $\sigma$²), then
\[aX + b \sim N(a\mu + b, a^2\sigma^2)\]

\textbf{Sum of Normals:} If X $\sim$ N($\mu$$_1$, $\sigma$$_1$²) and Y $\sim$ N($\mu$$_2$, $\sigma$$_2$²) independent, then
\[X + Y \sim N(\mu_1 + \mu_2, \sigma_1^2 + \sigma_2^2)\]
\end{important}

\begin{example}
\textbf{Distribution Selection Guide:}
- Binary outcome → Bernoulli
- Count of successes in n trials → Binomial
- Rare events/arrivals → Poisson
- Continuous, symmetric → Normal/Gaussian
\end{example}

\newpage

% ============================================================
% HYPOTHESIS TESTING
% ============================================================

\section{Hypothesis Testing}

\begin{important}{Hypothesis Testing Framework - VERY HIGH PROBABILITY}
\textbf{Step 1: State Hypotheses}
\begin{align*}
H_0&: \text{Null hypothesis (no effect, equality)} \\
H_1 \text{ or } H_a&: \text{Alternative hypothesis}
\end{align*}

\textbf{Types of Tests:}
- Two-tailed: H$_1$: $\mu$ $\neq$ $\mu$$_0$
- Right-tailed: H$_1$: $\mu$ > $\mu$$_0$
- Left-tailed: H$_1$: $\mu$ < $\mu$$_0$

\textbf{Step 2: Choose Significance Level}
\[\alpha = 0.05 \text{ (common), or } 0.01, 0.10\]

\textbf{Step 3: Calculate Test Statistic}

\textbf{Step 4: Find P-value or Critical Value}

\textbf{Step 5: Decision}
- If p-value < $\alpha$: Reject H$_0$
- If p-value $\geq$ $\alpha$: Fail to reject H$_0$
\end{important}

\begin{important}{Z-Test (Known $\sigma$) - HIGH PROBABILITY}
\textbf{Test Statistic:}
\[Z = \frac{\bar{x} - \mu_0}{\sigma/\sqrt{n}}\]

\textbf{Critical Values:}
- Two-tailed ($\alpha$=0.05): ±1.96
- One-tailed ($\alpha$=0.05): ±1.645
- Two-tailed ($\alpha$=0.01): ±2.576

\textbf{When to use:} Known population $\sigma$, large sample (n $\geq$ 30), or normal population.
\end{important}

\begin{important}{T-Test (Unknown $\sigma$) - HIGH PROBABILITY}
\textbf{Test Statistic:}
\[t = \frac{\bar{x} - \mu_0}{s/\sqrt{n}}\]

\textbf{Degrees of Freedom:} df = n - 1

\textbf{When to use:} Unknown $\sigma$, small sample, approximately normal data.

\textbf{Two-Sample T-Test (Equal Variances):}
\[t = \frac{\bar{x}_1 - \bar{x}_2}{s_p\sqrt{\frac{1}{n_1} + \frac{1}{n_2}}}\]

where pooled SD: \(s_p = \sqrt{\frac{(n_1-1)s_1^2 + (n_2-1)s_2^2}{n_1+n_2-2}}\)

df = n$_1$ + n$_2$ - 2
\end{important}

\begin{formula}{Type I and Type II Errors}
\begin{center}
\begin{tabular}{c|cc}
 & H$_0$ True & H$_0$ False \\
\hline
Reject H$_0$ & Type I Error ($\alpha$) & Correct (Power) \\
Fail to Reject H$_0$ & Correct (1-$\alpha$) & Type II Error ($\beta$) \\
\end{tabular}
\end{center}

\textbf{Type I Error:} False positive, $\alpha$ = P(Reject H$_0$ | H$_0$ true)

\textbf{Type II Error:} False negative, $\beta$ = P(Fail to reject H$_0$ | H$_0$ false)

\textbf{Power:} 1 - $\beta$ = P(Reject H$_0$ | H$_0$ false)
\end{formula}

\begin{formula}{Confidence Intervals}
\textbf{For Mean (known $\sigma$):}
\[\bar{x} \pm z_{\alpha/2}\frac{\sigma}{\sqrt{n}} = \left[\bar{x} - z_{\alpha/2}\frac{\sigma}{\sqrt{n}}, \bar{x} + z_{\alpha/2}\frac{\sigma}{\sqrt{n}}\right]\]

\textbf{For Mean (unknown $\sigma$):}
\[\bar{x} \pm t_{\alpha/2,n-1}\frac{s}{\sqrt{n}}\]

\textbf{For Proportion:}
\[\hat{p} \pm z_{\alpha/2}\sqrt{\frac{\hat{p}(1-\hat{p})}{n}}\]

\textbf{For Difference of Means (Independent Samples):}
\[(\bar{x}_1 - \bar{x}_2) \pm t_{\alpha/2,df}\sqrt{\frac{s_1^2}{n_1} + \frac{s_2^2}{n_2}}\]

\textbf{For Variance (Chi-Square):}
\[\left[\frac{(n-1)s^2}{\chi^2_{\alpha/2,n-1}}, \frac{(n-1)s^2}{\chi^2_{1-\alpha/2,n-1}}\right]\]

\textbf{Common Critical Values:}
- 90\% CI: $z_{0.05} = 1.645$, $t_{0.05,\infty} = 1.645$
- 95\% CI: $z_{0.025} = 1.96$, $t_{0.025,30} \approx 2.042$
- 99\% CI: $z_{0.005} = 2.576$, $t_{0.005,30} \approx 2.750$

\textbf{Interpretation:} $100(1-\alpha)\%$ of intervals constructed this way contain true parameter.
\end{formula}

\begin{important}{Effect Sizes - HIGH PROBABILITY}
\textbf{Cohen's d (Standardized Mean Difference):}
\[d = \frac{\bar{x}_1 - \bar{x}_2}{s_{\text{pooled}}} \quad \text{where} \quad s_{\text{pooled}} = \sqrt{\frac{(n_1-1)s_1^2 + (n_2-1)s_2^2}{n_1+n_2-2}}\]
Interpretation: $|d| = 0.2$ (small), $0.5$ (medium), $0.8$ (large)

\textbf{Cohen's d (Single Sample):}
\[d = \frac{\bar{x} - \mu_0}{s}\]

\textbf{Hedges' g (Corrected for Small Samples):}
\[g = d \times \left(1 - \frac{3}{4(n_1 + n_2) - 9}\right)\]

\textbf{Eta-Squared (Effect Size for ANOVA):}
\[\eta^2 = \frac{\text{SS}_{\text{between}}}{\text{SS}_{\text{total}}} = \frac{\text{SS}_B}{\text{SS}_B + \text{SS}_W}\]

\textbf{Omega-Squared (Less Biased):}
\[\omega^2 = \frac{\text{SS}_B - (k-1)\text{MS}_W}{\text{SS}_{\text{total}} + \text{MS}_W}\]
\end{important}

\begin{example}
\textbf{Quick Decision Rules:}
- p-value < 0.05 → Reject H$_0$ (significant result)
- |Z| > 1.96 → Reject H$_0$ at 5\% level (two-tailed)
- Confidence interval not containing H$_0$ value → Reject H$_0$
\end{example}

\begin{important}{Chi-Square Test - HIGH PROBABILITY}
\textbf{Goodness of Fit Test:}
\[\chi^2 = \sum_{i=1}^{k}\frac{(O_i - E_{i})^2}{E_i}\]

where:
- $O_{i}$ = observed frequency
- $E_{i}$ = expected frequency
- k = number of categories

df = k - 1

\textbf{Test of Independence (Contingency Table):}
\[\chi^2 = \sum_{i=1}^{r}\sum_{j=1}^{c}\frac{(O_{ij} - E_{ij})^2}{E_{ij}}\]

\textbf{Expected Frequency:}
\[E_{ij} = \frac{(\text{Row}_i \text{ Total}) \times (\text{Col}_j \text{ Total})}{\text{Grand Total}}\]

df = (r-1)(c-1)

\textbf{Assumption:} All expected frequencies $\geq$ 5

\textbf{Effect Size - Cramér's V:}
\[V = \sqrt{\frac{\chi^2}{n \times \min(r-1, c-1)}}\]

Interpretation: 0.1 (small), 0.3 (medium), 0.5 (large)
\end{important}

\begin{formula}{ANOVA (Analysis of Variance) - MEDIUM PROBABILITY}
\textbf{One-Way ANOVA F-statistic:}
\[F = \frac{MSB}{MSW} = \frac{SSB/(k-1)}{SSW/(N-k)}\]

where:
- k = number of groups
- N = total sample size
- SSB = Sum of Squares Between groups
- SSW = Sum of Squares Within groups

\textbf{Sum of Squares:}
\[SST = \sum_{i=1}^{k}\sum_{j=1}^{n_i}(x_{ij} - \bar{x})^2\]

\[SSB = \sum_{i=1}^{k}n_i(\bar{x}_i - \bar{x})^2\]

\[SSW = \sum_{i=1}^{k}\sum_{j=1}^{n_i}(x_{ij} - \bar{x}_i)^2\]

\textbf{Relationship:} SST = SSB + SSW

\textbf{Degrees of Freedom:}
- df(Between) = k - 1
- df(Within) = N - k
- df(Total) = N - 1

\textbf{Hypothesis:}
- H$_0$: $\mu$$_1$ = $\mu$$_2$ = ... = $\mu$$_k$
- H$_1$: At least one mean differs

\textbf{Effect Size - Eta-squared:}
\[\eta^2 = \frac{SSB}{SST}\]
\end{formula}

\begin{formula}{Paired T-Test - MEDIUM PROBABILITY}
\textbf{For paired/dependent samples:}
\[t = \frac{\bar{d}}{s_d/\sqrt{n}}\]

where:
- $d_{i}$ = x_{i1} - x_{i2} (difference for pair i)
- \(\bar{d}\) = mean of differences
- $s_{d}$ = standard deviation of differences
- df = n - 1

\textbf{When to use:} Before-after measurements, matched pairs
\end{formula}

\begin{formula}{Bootstrap and Resampling - LOW PROBABILITY}
\textbf{Bootstrap Method:}
1. Draw n samples with replacement from original data
2. Calculate statistic of interest (mean, median, etc.)
3. Repeat B times (typically B = 1000-10000)
4. Use distribution of bootstrap statistics for inference

\textbf{Bootstrap Confidence Interval (Percentile method):}
\[CI = [$\theta$_{\alpha/2}^*, $\theta$_{1-\alpha/2}^*]\]

where $\theta$* are bootstrap estimates.

\textbf{When to use:} Unknown distribution, small samples, complex statistics
\end{formula}

\begin{formula}{Cross-Validation}
\textbf{K-Fold Cross-Validation:}
1. Split data into K folds
2. Train on K-1 folds, test on remaining fold
3. Repeat K times
4. Average performance metrics

\textbf{Cross-Validation Error:}
\[CV_{(K)} = \frac{1}{K}\sum_{i=1}^{K} \text{Error}_i\]

\textbf{Leave-One-Out CV (LOOCV):} K = n (special case)

\textbf{When to use:} Model selection, hyperparameter tuning, avoiding overfitting
\end{formula}

\begin{important}{Min-Max Normalization - HIGH PROBABILITY}
\textbf{Formula:} Scales data to [0, 1] range
\[x_{\text{norm}} = \frac{x - x_{\min}}{x_{\max} - x_{\min}}\]

\textbf{General range [a, b]:}
\[x_{\text{norm}} = a + \frac{(x - x_{\min})(b - a)}{x_{\max} - x_{\min}}\]

\textbf{When to use:} Neural networks, image processing, distance-based algorithms (KNN, K-Means).

\textbf{Pros:} Preserves relationships, bounded range.

\textbf{Cons:} Sensitive to outliers.
\end{important}

\begin{important}{Z-Score Normalization (Standardization) - HIGH PROBABILITY}
\textbf{Formula:}
\[z = \frac{x - \mu}{\sigma} = \frac{x - \bar{x}}{s}\]

\textbf{Properties:}
- Mean = 0
- Standard deviation = 1
- Unbounded range

\textbf{When to use:} When data should have mean 0 and SD 1, linear regression, PCA, algorithms assuming normal distribution.

\textbf{Pros:} Not bounded, handles outliers better than min-max.

\textbf{Cons:} Doesn't preserve exact relationships.
\end{important}

\begin{formula}{One-Hot Encoding - MEDIUM PROBABILITY}
\textbf{For categorical variable with k categories:}

Create k binary columns, each representing one category.

\textbf{Example:} Color $\in$ \{Red, Green, Blue\}

\begin{center}
\begin{tabular}{c|ccc}
Original & Red & Green & Blue \\
\hline
Red & 1 & 0 & 0 \\
Green & 0 & 1 & 0 \\
Blue & 0 & 0 & 1 \\
\end{tabular}
\end{center}

\textbf{Formula:}
\[x_{i}^{(j)} = \begin{cases}
1 & \text{if category } = j \\
0 & \text{otherwise}
\end{cases}\]

\textbf{When to use:} Nominal categorical variables with no ordering.

\textbf{Note:} For linear models, use k-1 columns to avoid multicollinearity (dummy variable trap).
\end{formula}

\begin{formula}{Label Encoding}
\textbf{For ordinal categorical variables:}

Assign integers 0, 1, 2, ... to categories based on order.

\textbf{Example:} Size $\in$ \{Small, Medium, Large\}
- Small → 0
- Medium → 1
- Large → 2

\textbf{When to use:} Ordinal data with natural ordering, tree-based models.
\end{formula}

\begin{formula}{Robust Scaling - MEDIUM PROBABILITY}
\textbf{Formula:} Uses median and IQR to be robust to outliers
\[x_{\text{scaled}} = \frac{x - \text{Median}(x)}{IQR}\]

where IQR = Q$_3$ - Q$_1$

\textbf{When to use:} Data with outliers, non-normal distributions

\textbf{Advantages:} Not affected by extreme values, preserves shape
\end{formula}

\begin{formula}{Log Transformation}
\textbf{For skewed data:}
\[x_{\text{transformed}} = \log(x)\]

or for values close to zero:
\[x_{\text{transformed}} = \log(x + 1)\]

\textbf{Box-Cox Transformation:}
\[x^{(\lambda)} = \begin{cases}
\frac{x^{\lambda} - 1}{\lambda} & \text{if } \lambda \neq 0 \\
\log(x) & \text{if } \lambda = 0
\end{cases}\]

\textbf{When to use:} Highly skewed data, making data more normal, stabilizing variance
\end{formula}

\textbf{Least Squares Estimates:}
\[\beta_1 = \frac{\sum_{i=1}^{n}(x_{i} - \bar{x})(y_{i} - \bar{y})}{\sum_{i=1}^{n}(x_{i} - \bar{x})^2} = \frac{S_{xy}}{S_{xx}}\]

\[\beta_0 = \bar{y} - \beta_1\bar{x}\]

\textbf{Alternative formula for slope:}
\[\beta_1 = r \cdot \frac{s_y}{s_x}\]

where r is correlation coefficient.

\textbf{Prediction:}
\[\hat{y} = \beta_0 + \beta_1 x\]
\end{important}

\begin{important}{Multiple Linear Regression - HIGH PROBABILITY}
\textbf{Model:}
\[y = \beta_0 + \beta_1 x_1 + \beta_2 x_2 + \cdots + \beta_p x_p + \epsilon\]

\textbf{Matrix Form:}
\[\mathbf{y} = \mathbf{X}\boldsymbol{\beta} + \boldsymbol{\epsilon}\]

\textbf{OLS Solution:}
\[\hat{\boldsymbol{\beta}} = (\mathbf{X}^T\mathbf{X})^{-1}\mathbf{X}^T\mathbf{y}\]

\textbf{Prediction:}
\[\hat{\mathbf{y}} = \mathbf{X}\hat{\boldsymbol{\beta}}\]
\end{important}

\begin{important}{Regression Metrics - VERY HIGH PROBABILITY}
\textbf{Residuals:}
\[e_i = y_{i} - \hat{y}_i\]

\textbf{Sum of Squared Errors (SSE):}
\[SSE = \sum_{i=1}^{n}(y_{i} - \hat{y}_i)^2 = \sum_{i=1}^{n}e_i^2\]

\textbf{Total Sum of Squares (SST):}
\[SST = \sum_{i=1}^{n}(y_{i} - \bar{y})^2\]

\textbf{Regression Sum of Squares (SSR):}
\[SSR = \sum_{i=1}^{n}(\hat{y}_i - \bar{y})^2\]

\textbf{Relationship:} SST = SSR + SSE

\textbf{R-squared (Coefficient of Determination):}
\[R^2 = \frac{SSR}{SST} = 1 - \frac{SSE}{SST}\]

Interpretation: Proportion of variance in y explained by x.

\textbf{Adjusted R-squared:}
\[R^2_{\text{adj}} = 1 - \frac{SSE/(n-p-1)}{SST/(n-1)}\]

Penalizes adding unnecessary predictors.

\textbf{Mean Squared Error (MSE):}
\[MSE = \frac{1}{n}\sum_{i=1}^{n}(y_{i} - \hat{y}_i)^2 = \frac{SSE}{n}\]

\textbf{Root Mean Squared Error (RMSE):}
\[RMSE = \sqrt{MSE} = \sqrt{\frac{SSE}{n}}\]

\textbf{Mean Absolute Error (MAE):}
\[MAE = \frac{1}{n}\sum_{i=1}^{n}|y_{i} - \hat{y}_i|\]
\end{important}

\subsection{Logistic Regression}

\begin{important}{Logistic Regression - HIGH PROBABILITY}
\textbf{Model:} For binary outcome (0 or 1)
\[P(Y=1|X) = \frac{1}{1 + e^{-(\beta_0 + \beta_1 x_1 + \cdots + \beta_p x_p)}}\]

\textbf{Logit (Log-Odds):}
\[\log\left(\frac{p}{1-p}\right) = \beta_0 + \beta_1 x_1 + \cdots + \beta_p x_p\]

\textbf{Odds:}
\[\text{Odds} = \frac{p}{1-p} = e^{\beta_0 + \beta_1 x_1 + \cdots + \beta_p x_p}\]

\textbf{Sigmoid Function:}
\[\sigma(z) = \frac{1}{1 + e^{-z}}\]

\textbf{Decision Rule:}
\[\hat{y} = \begin{cases}
1 & \text{if } P(Y=1|X) \geq 0.5 \\
0 & \text{if } P(Y=1|X) < 0.5
\end{cases}\]

\textbf{When to use:} Binary classification problems.
\end{important}

\subsection{Polynomial Regression}

\begin{formula}{Polynomial Regression - MEDIUM PROBABILITY}
\textbf{Degree 2 (Quadratic):}
\[y = \beta_0 + \beta_1 x + \beta_2 x^2 + \epsilon\]

\textbf{Degree p:}
\[y = \beta_0 + \beta_1 x + \beta_2 x^2 + \cdots + \beta_p x^p + \epsilon\]

\textbf{Note:} Still linear in parameters $\beta$, so use OLS.

\textbf{When to use:} Non-linear relationships, curved patterns in data.

\textbf{Warning:} High degrees can cause overfitting.
\end{formula}

\subsection{Ridge Regression (L2 Regularization)}

\begin{important}{Ridge Regression - MEDIUM PROBABILITY}
\textbf{Objective:} Minimize
\[SSE + \lambda\sum_{j=1}^{p}\beta_j^2\]

\textbf{Solution:}
\[\hat{\boldsymbol{\beta}}_{\text{ridge}} = (\mathbf{X}^T\mathbf{X} + \lambda\mathbf{I})^{-1}\mathbf{X}^T\mathbf{y}\]

where $\lambda$ $\geq$ 0 is regularization parameter.

\textbf{Effect:}
- $\lambda$ = 0: Ordinary least squares
- $\lambda$ > 0: Shrinks coefficients toward zero
- $\lambda$ → $\infty$: All coefficients → 0

\textbf{When to use:} Multicollinearity, many predictors, prevent overfitting.

\textbf{Note:} Standardize features before applying ridge regression.
\end{important}

\begin{formula}{LASSO Regression (L1 Regularization)}
\textbf{Objective:} Minimize
\[SSE + \lambda\sum_{j=1}^{p}|\beta_j|\]

\textbf{Effect:} Can shrink some coefficients exactly to zero (feature selection).

\textbf{Comparison:}
- Ridge: Shrinks but keeps all features
- LASSO: Can eliminate features (sparse solutions)

\textbf{Elastic Net:} Combines L1 and L2
\[SSE + \lambda_1\sum|\beta_j| + \lambda_2\sum\beta_j^2\]
\end{formula}

\begin{formula}{Logistic Regression Cost Function}
\textbf{Log-Likelihood (Cost Function):}
\[L(\beta) = -\sum_{i=1}^{n}\left[y_{i}\log(p_i) + (1-y_{i})\log(1-p_i)\right]\]

where \(p_i = \frac{1}{1 + e^{-(\beta_0 + \beta_1 x_{i})}}\)

\textbf{Optimization:} Use gradient descent or Newton-Raphson

\textbf{Odds Ratio:} For one unit increase in x
\[OR = e^{\beta_1}\]

Interpretation: Multiplicative change in odds
\end{formula}

\begin{important}{Residual Analysis - MEDIUM PROBABILITY}
\textbf{Residual:}
\[e_i = y_{i} - \hat{y}_i\]

\textbf{Standardized Residual:}
\[r_i = \frac{e_i}{s\sqrt{1 - h_{ii}}}\]

where h_{ii} is leverage (diagonal of hat matrix)

\textbf{Cook's Distance (influence):}
\[D_i = \frac{r_i^2}{p}\frac{h_{ii}}{1-h_{ii}}\]

\textbf{Rules of thumb:}
- |r_i| > 2: Potential outlier
- D_i > 4/n: Influential point

\textbf{Durbin-Watson Test (autocorrelation):}
\[DW = \frac{\sum_{i=2}^{n}(e_i - e_{i-1})^2}{\sum_{i=1}^{n}e_i^2}\]

Range: 0 $\leq$ DW $\leq$ 4
- DW $\approx$ 2: No autocorrelation
- DW < 2: Positive autocorrelation
- DW > 2: Negative autocorrelation
\end{important}

% ============================================================
% DECISION TREES
% ============================================================

\section{Decision Trees}

\begin{important}{Entropy - HIGH PROBABILITY}
\textbf{Definition:} Measure of impurity/disorder
\[H(S) = -\sum_{i=1}^{c} p_i \log_2(p_i)\]

where:
- c = number of classes
- p_i = proportion of class i in set S

\textbf{Properties:}
- H = 0: Pure (all same class)
- H = 1: Maximum impurity (binary, 50-50 split)
- Higher entropy = more disorder

\textbf{For binary classification:}
\[H = -p\log_2(p) - (1-p)\log_2(1-p)\]
\end{important}

\begin{important}{Information Gain - HIGH PROBABILITY}
\textbf{Formula:}
\[IG(S, A) = H(S) - \sum_{v \in \text{Values}(A)} \frac{|S_v|}{|S|} H(S_v)\]

where:
- S = parent set
- A = attribute/feature
- S_v = subset of S where attribute A has value v

\textbf{Interpretation:} Reduction in entropy after split on attribute A.

\textbf{Decision Rule:} Choose attribute with maximum information gain.
\end{important}

\begin{important}{Gini Index (Gini Impurity) - HIGH PROBABILITY}
\textbf{Formula:}
\[Gini(S) = 1 - \sum_{i=1}^{c} p_i^2\]

\textbf{For binary classification:}
\[Gini = 1 - p^2 - (1-p)^2 = 2p(1-p)\]

\textbf{Properties:}
- Gini = 0: Pure (all same class)
- Gini = 0.5: Maximum impurity (binary, 50-50 split)
- Lower Gini = less impurity

\textbf{Gini Gain:}
\[GiniGain(S, A) = Gini(S) - \sum_{v} \frac{|S_v|}{|S|} Gini(S_v)\]

\textbf{Note:} CART (Classification and Regression Trees) uses Gini Index.
\end{important}

\begin{example}
\textbf{Entropy vs Gini:}
- Both measure impurity
- Entropy: Logarithmic, ranges [0, 1] for binary
- Gini: Quadratic, ranges [0, 0.5] for binary
- Gini is faster to compute (no logarithms)
- Results usually similar in practice
\end{example}

\begin{formula}{Decision Tree Pruning}
\textbf{Cost-Complexity Pruning ($\alpha$ parameter):}
\[\text{Cost} = \text{Error}(T) + \alpha \times |T|\]

where:
- Error(T) = misclassification error of tree T
- |T| = number of terminal nodes
- $\alpha$ = complexity parameter

\textbf{Goal:} Find subtree that minimizes cost

\textbf{Methods:}
- Reduced Error Pruning
- Cost-Complexity Pruning (weakest link)
- Pessimistic Error Pruning

\textbf{When to prune:} Prevent overfitting, improve generalization
\end{formula}

\begin{formula}{Random Forest}
\textbf{Ensemble Method:} Combines multiple decision trees

\textbf{Algorithm:}
1. Create B bootstrap samples
2. Build tree on each sample
3. For each split, consider m random features (m < p)
4. Aggregate predictions (majority vote or average)

\textbf{Out-of-Bag (OOB) Error:}
\[\text{OOB Error} = \frac{1}{n}\sum_{i=1}^{n}I(y_{i} \neq \hat{y}_i^{OOB})\]

\textbf{Feature Importance:}
Based on decrease in impurity when feature is used for splitting

\textbf{Typical values:} m = √p (classification), m = p/3 (regression)
\end{formula}

% ============================================================
% CONFUSION MATRIX AND METRICS
% ============================================================

\section{Confusion Matrix and Classification Metrics}

\begin{important}{Confusion Matrix - VERY HIGH PROBABILITY}
\textbf{For Binary Classification:}

\begin{center}
\begin{tabular}{cc|cc}
& & \multicolumn{2}{c}{Predicted} \\
& & Positive & Negative \\
\hline
\multirow{2}{*}{Actual} & Positive & TP & FN \\
& Negative & FP & TN \\
\end{tabular}
\end{center}

where:
- TP = True Positives (correctly predicted positive)
- TN = True Negatives (correctly predicted negative)
- FP = False Positives (Type I error)
- FN = False Negatives (Type II error)
\end{important}

\begin{important}{Classification Metrics - VERY HIGH PROBABILITY}
\textbf{Accuracy:}
\[Accuracy = \frac{TP + TN}{TP + TN + FP + FN}\]

\textbf{Precision (Positive Predictive Value):}
\[Precision = \frac{TP}{TP + FP}\]

Interpretation: Of all predicted positives, what fraction are actually positive?

\textbf{Recall (Sensitivity, True Positive Rate):}
\[Recall = \frac{TP}{TP + FN}\]

Interpretation: Of all actual positives, what fraction did we predict?

\textbf{Specificity (True Negative Rate):}
\[Specificity = \frac{TN}{TN + FP}\]

\textbf{False Positive Rate:}
\[FPR = \frac{FP}{FP + TN} = 1 - Specificity\]

\textbf{F1-Score (Harmonic Mean of Precision and Recall):}
\[F1 = 2 \cdot \frac{Precision \times Recall}{Precision + Recall} = \frac{2TP}{2TP + FP + FN}\]

\textbf{F-beta Score:}
\[F_\beta = (1 + \beta^2) \cdot \frac{Precision \times Recall}{\beta^2 \cdot Precision + Recall}\]

- $\beta$ > 1: Emphasizes recall
- $\beta$ < 1: Emphasizes precision
\end{important}

\begin{example}
\textbf{Metric Selection Guide:}
- Balanced classes → Accuracy
- Imbalanced classes → F1-Score, Precision/Recall
- Cost of FP high (spam detection) → Precision
- Cost of FN high (disease screening) → Recall
- Trade-off between FP and FN → F1-Score
\end{example}

\begin{formula}{ROC Curve and AUC - MEDIUM PROBABILITY}
\textbf{Receiver Operating Characteristic (ROC):}

Plot of True Positive Rate vs False Positive Rate at different thresholds.

\textbf{TPR (Sensitivity):}
\[TPR = \frac{TP}{TP + FN}\]

\textbf{FPR:}
\[FPR = \frac{FP}{FP + TN}\]

\textbf{Area Under Curve (AUC):}
- AUC = 1.0: Perfect classifier
- AUC = 0.5: Random classifier
- AUC > 0.7: Generally acceptable
- AUC > 0.8: Good model
- AUC > 0.9: Excellent model

\textbf{When to use:} Comparing classifier performance, threshold selection
\end{formula}

\begin{formula}{Matthews Correlation Coefficient (MCC)}
\textbf{For imbalanced datasets:}
\[MCC = \frac{TP \times TN - FP \times FN}{\sqrt{(TP+FP)(TP+FN)(TN+FP)(TN+FN)}}\]

\textbf{Range:} -1 $\leq$ MCC $\leq$ 1
- MCC = 1: Perfect prediction
- MCC = 0: Random prediction
- MCC = -1: Total disagreement

\textbf{Advantage:} Balanced measure even for imbalanced classes
\end{formula}

\begin{formula}{Cohen's Kappa - LOW PROBABILITY}
\textbf{Agreement measure:}
\[\kappa = \frac{p_o - p_e}{1 - p_e}\]

where:
- p_o = observed agreement
- p_e = expected agreement by chance

\textbf{Interpretation:}
- κ < 0: No agreement
- 0.01-0.20: Slight agreement
- 0.21-0.40: Fair agreement
- 0.41-0.60: Moderate agreement
- 0.61-0.80: Substantial agreement
- 0.81-1.00: Almost perfect agreement
\end{formula}

% ============================================================
% NAIVE BAYES
% ============================================================

\section{Naive Bayes Classifier}

\begin{important}{Bayes' Theorem - VERY HIGH PROBABILITY}
\textbf{Formula:}
\[P(A|B) = \frac{P(B|A) \cdot P(A)}{P(B)}\]

\textbf{For Classification:}
\[P(C_k|X) = \frac{P(X|C_k) \cdot P(C_k)}{P(X)}\]

where:
- C_k = class k
- X = feature vector
- P(C_k) = prior probability of class k
- P(X|C_k) = likelihood
- P(C_k|X) = posterior probability
\end{important}

\begin{important}{Naive Bayes Classification - HIGH PROBABILITY}
\textbf{Assumption:} Features are conditionally independent given class.

\[P(C_k|x_1, x_2, \ldots, x_n) \propto P(C_k) \prod_{i=1}^{n} P(x_{i}|C_k)\]

\textbf{Decision Rule:} Assign to class with maximum posterior probability
\[\hat{y} = \arg\max_{k} P(C_k) \prod_{i=1}^{n} P(x_{i}|C_k)\]

\textbf{Log Form (to avoid underflow):}
\[\hat{y} = \arg\max_{k} \left[\log P(C_k) + \sum_{i=1}^{n} \log P(x_{i}|C_k)\right]\]
\end{important}

\begin{formula}{Naive Bayes for Continuous Features (Gaussian)}
\textbf{Likelihood:} Assume features follow normal distribution
\[P(x_{i}|C_k) = \frac{1}{\sqrt{2\pi\sigma_{k,i}^2}} \exp\left(-\frac{(x_{i} - \mu_{k,i})^2}{2\sigma_{k,i}^2}\right)\]

where $\mu$_{k,i} and $\sigma$_{k,i} are mean and SD of feature i in class k.
\end{formula}

\begin{formula}{Naive Bayes for Discrete Features (Multinomial)}
\textbf{Likelihood:}
\[P(x_{i}|C_k) = \frac{\text{count}(x_{i}, C_k) + \alpha}{\text{count}(C_k) + \alpha \cdot |V|}\]

where:
- $\alpha$ = smoothing parameter (Laplace smoothing, typically $\alpha$ = 1)
- |V| = vocabulary size (number of unique values)

\textbf{Laplace Smoothing:} Prevents zero probabilities

\textbf{Add-k Smoothing:} Generalization with $\alpha$ = k
\end{formula}

\begin{formula}{Naive Bayes Model Selection}
\textbf{Gaussian Naive Bayes:}
- Continuous features
- Assumes normal distribution per class

\textbf{Multinomial Naive Bayes:}
- Discrete count data
- Text classification (word counts)

\textbf{Bernoulli Naive Bayes:}
- Binary features (present/absent)
- Document classification (word presence)

\textbf{Advantages:}
- Fast training and prediction
- Works well with small datasets
- Handles high-dimensional data
- Probabilistic output

\textbf{Disadvantages:}
- Independence assumption rarely holds
- Can't capture feature interactions
\end{formula}

% ============================================================
% K-NEAREST NEIGHBORS (KNN)
% ============================================================

\section{K-Nearest Neighbors (KNN)}

\begin{important}{KNN Algorithm - HIGH PROBABILITY}
\textbf{Steps:}
1. Choose number of neighbors k
2. Calculate distance to all training points
3. Select k nearest neighbors
4. For classification: Majority vote
5. For regression: Average of k neighbors

\textbf{Classification:}
\[\hat{y} = \text{mode}\{y_1, y_2, \ldots, y_k\}\]

\textbf{Regression:}
\[\hat{y} = \frac{1}{k}\sum_{i=1}^{k} y_{i}\]
\end{important}

\begin{important}{Distance Metrics - HIGH PROBABILITY}
\textbf{Euclidean Distance (L2):}
\[d(x, x') = \sqrt{\sum_{i=1}^{n}(x_{i} - x'_i)^2}\]

\textbf{Manhattan Distance (L1):}
\[d(x, x') = \sum_{i=1}^{n}|x_{i} - x'_i|\]

\textbf{Minkowski Distance (generalized):}
\[d(x, x') = \left(\sum_{i=1}^{n}|x_{i} - x'_i|^p\right)^{1/p}\]

- p = 1: Manhattan
- p = 2: Euclidean
- p → $\infty$: Chebyshev (max difference)

\textbf{Cosine Similarity:}
\[\text{similarity} = \frac{\mathbf{x} \cdot \mathbf{x}'}{\|\mathbf{x}\| \|\mathbf{x}'\|} = \frac{\sum x_{i} x'_i}{\sqrt{\sum x_{i}^2}\sqrt{\sum x_{i}'^2}}\]

Cosine Distance: d = 1 - similarity
\end{important}

\begin{formula}{Weighted KNN}
\textbf{Weights inversely proportional to distance:}
\[w_i = \frac{1}{d_i}\]

\textbf{Weighted Average:}
\[\hat{y} = \frac{\sum_{i=1}^{k} w_i y_{i}}{\sum_{i=1}^{k} w_i}\]

\textbf{When to use:} Give more importance to closer neighbors.
\end{formula}

\begin{example}
\textbf{Choosing k:}
- k too small: Overfitting, sensitive to noise
- k too large: Underfitting, over-smoothing
- Rule of thumb: k = √n (n = training size)
- Always use odd k for binary classification (avoid ties)
- Use cross-validation to find optimal k
\end{example}

\begin{formula}{KNN for Imbalanced Data}
\textbf{Weighted Voting:}

Give more weight to closer neighbors and adjust for class imbalance.

\textbf{Weight:}
\[w_i = \frac{1}{d_i} \times \frac{n}{n_c}\]

where:
- $d_{i}$ = distance to neighbor i
- n = total samples
- n_c = samples in class c

\textbf{Distance-Weighted KNN:} Closer neighbors have more influence

\textbf{Class-Balanced KNN:} Adjusts for class distribution
\end{formula}

\begin{formula}{Curse of Dimensionality}
\textbf{Problem:} As dimensions increase, distance becomes less meaningful.

\textbf{Effect on KNN:}
- All points become roughly equidistant
- Nearest neighbors aren't truly "near"
- Performance degrades

\textbf{Solutions:}
- Dimensionality reduction (PCA, feature selection)
- Feature engineering
- Use appropriate distance metrics
- Increase sample size exponentially with dimensions

\textbf{Rule of thumb:} Keep dimensions p << n (number of samples)
\end{formula}

% ============================================================
% K-MEANS CLUSTERING
% ============================================================

\section{K-Means Clustering}

\begin{important}{K-Means Algorithm - HIGH PROBABILITY}
\textbf{Objective:} Minimize within-cluster sum of squares (WCSS)
\[J = \sum_{i=1}^{k}\sum_{x \in C_i} \|x - \mu_i\|^2\]

where:
- k = number of clusters
- C_i = cluster i
- $\mu$_i = centroid of cluster i

\textbf{Algorithm Steps:}
1. Initialize k centroids randomly
2. \textbf{Assignment step:} Assign each point to nearest centroid
   \[C_i = \{x : \|x - \mu_i\| \leq \|x - \mu_j\| \text{ for all } j\}\]
3. \textbf{Update step:} Recalculate centroids
   \[\mu_i = \frac{1}{|C_i|}\sum_{x \in C_i} x\]
4. Repeat steps 2-3 until convergence
\end{important}

\begin{important}{Within-Cluster Sum of Squares (WCSS) - HIGH PROBABILITY}
\textbf{Formula:}
\[WCSS = \sum_{i=1}^{k}\sum_{x \in C_i} \|x - \mu_i\|^2\]

\textbf{Inertia:} Another name for WCSS.

\textbf{Elbow Method:} Plot WCSS vs k, choose k at "elbow" where rate of decrease slows.
\end{important}

\begin{formula}{Silhouette Score - MEDIUM PROBABILITY}
\textbf{For point i:}
\[s(i) = \frac{b(i) - a(i)}{\max\{a(i), b(i)\}}\]

where:
- a(i) = average distance to points in same cluster
- b(i) = average distance to points in nearest other cluster

\textbf{Range:} -1 $\leq$ s(i) $\leq$ 1
- s(i) $\approx$ 1: Well clustered
- s(i) $\approx$ 0: On cluster boundary
- s(i) < 0: Wrong cluster

\textbf{Overall Silhouette Score:}
\[s = \frac{1}{n}\sum_{i=1}^{n} s(i)\]
\end{formula}

\begin{example}
\textbf{K-Means Tips:}
- Sensitive to initial centroids (use k-means++)
- Assumes spherical clusters of similar size
- Requires pre-specifying k
- Scale/normalize features before clustering
- Run multiple times with different initializations
\end{example}

\begin{formula}{K-Means++ Initialization}
\textbf{Algorithm:} Better initialization than random

1. Choose first centroid randomly
2. For each data point, compute distance D(x) to nearest centroid
3. Choose next centroid with probability $\propto$ D(x)²
4. Repeat until k centroids chosen

\textbf{Advantage:} Faster convergence, better final clusters
\end{formula}

\begin{formula}{Hierarchical Clustering}
\textbf{Linkage Methods:}

\textbf{Single Linkage:}
\[d(A, B) = \min_{a \in A, b \in B} d(a, b)\]

\textbf{Complete Linkage:}
\[d(A, B) = \max_{a \in A, b \in B} d(a, b)\]

\textbf{Average Linkage:}
\[d(A, B) = \frac{1}{|A||B|}\sum_{a \in A}\sum_{b \in B} d(a, b)\]

\textbf{Ward's Method:} Minimizes within-cluster variance
\[d(A, B) = \frac{|A||B|}{|A|+|B|} \|\mu_A - \mu_B\|^2\]

\textbf{Dendrogram:} Tree diagram showing cluster hierarchy

\textbf{When to use:} Unknown number of clusters, visualizing relationships
\end{formula}

\begin{formula}{Gap Statistic for Choosing k}
\textbf{Formula:}
\[Gap(k) = E[\log(W_{k})] - \log(W_{k})\]

where:
- W_{k} = within-cluster sum of squares for k clusters
- E[log(W_{k})] = expected value under null reference distribution

\textbf{Optimal k:} Choose smallest k where
\[Gap(k) \geq Gap(k+1) - s_{k+1}\]

where s_{k+1} is standard deviation.

\textbf{Interpretation:} Compares clustering to random data
\end{formula}

% ============================================================
% PRINCIPAL COMPONENT ANALYSIS (PCA)
% ============================================================

\section{Principal Component Analysis (PCA)}

\begin{important}{PCA Overview - MEDIUM PROBABILITY}
\textbf{Goal:} Reduce dimensionality by finding directions of maximum variance.

\textbf{Steps:}
1. Standardize the data (mean 0, variance 1)
2. Compute covariance matrix
3. Calculate eigenvalues and eigenvectors
4. Sort eigenvectors by eigenvalues (descending)
5. Select top k eigenvectors (principal components)
6. Transform data to new space
\end{important}

\begin{important}{Covariance Matrix - HIGH PROBABILITY}
\textbf{For standardized data X (n $\times$ p):}
\[\mathbf{C} = \frac{1}{n-1}\mathbf{X}^T\mathbf{X}\]

\textbf{Element (i,j):}
\[C_{ij} = \frac{1}{n-1}\sum_{k=1}^{n}(x_{ki} - \bar{x}_i)(x_{kj} - \bar{x}_j)\]

\textbf{Diagonal:} Variances of features

\textbf{Off-diagonal:} Covariances between features
\end{important}

\begin{important}{Eigenvalues and Eigenvectors - MEDIUM PROBABILITY}
\textbf{Eigenvalue equation:}
\[\mathbf{C}\mathbf{v} = \lambda\mathbf{v}\]

where:
- C = covariance matrix
- v = eigenvector (principal component direction)
- $\lambda$ = eigenvalue (variance along that direction)

\textbf{Properties:}
- Each eigenvalue corresponds to one principal component
- Larger eigenvalue = more variance explained
- Eigenvectors are orthogonal (perpendicular)
\end{important}

\begin{important}{Variance Explained - HIGH PROBABILITY}
\textbf{Total Variance:}
\[\text{Total Var} = \sum_{i=1}^{p} \lambda_i\]

\textbf{Variance Explained by PC i:}
\[\text{Var}_i = \frac{\lambda_i}{\sum_{j=1}^{p} \lambda_j}\]

\textbf{Cumulative Variance:}
\[\text{Cumulative Var}_k = \frac{\sum_{i=1}^{k} \lambda_i}{\sum_{j=1}^{p} \lambda_j}\]

\textbf{Rule of thumb:} Keep components explaining 80-95\% of variance.
\end{important}

\begin{formula}{PCA Transformation}
\textbf{Transform data to PC space:}
\[\mathbf{Z} = \mathbf{X}\mathbf{W}\]

where:
- X = original data (n $\times$ p)
- W = matrix of k eigenvectors (p $\times$ k)
- Z = transformed data (n $\times$ k)

\textbf{Reconstruction (approximate):}
\[\mathbf{X}_{\text{approx}} = \mathbf{Z}\mathbf{W}^T\]
\end{formula}

\begin{example}
\textbf{PCA Use Cases:}
- Dimensionality reduction (visualization, speed)
- Feature extraction
- Noise filtering
- Removing multicollinearity
- Data compression

\textbf{Note:} PCA is sensitive to scaling - always standardize first!
\end{example}

\begin{formula}{PCA Reconstruction Error}
\textbf{Mean Squared Reconstruction Error:}
\[\text{MSE} = \frac{1}{n}\|\mathbf{X} - \mathbf{X}_{\text{approx}}\|^2_F\]

where \(\|\cdot\|_F\) is Frobenius norm.

\textbf{Relation to Eigenvalues:}
\[\text{MSE} = \sum_{i=k+1}^{p} \lambda_i\]

\textbf{Interpretation:} Sum of discarded eigenvalues
\end{formula}

\begin{important}{Kernel PCA - LOW PROBABILITY}
\textbf{For non-linear relationships:}

Apply kernel trick to capture non-linear structure.

\textbf{Common Kernels:}

\textbf{RBF (Gaussian):}
\[K(x, y) = \exp\left(-\frac{\|x - y\|^2}{2\sigma^2}\right)\]

\textbf{Polynomial:}
\[K(x, y) = (x^T y + c)^d\]

\textbf{Sigmoid:}
\[K(x, y) = \tanh(\alpha x^T y + c)\]

\textbf{When to use:} Data with non-linear relationships, complex manifolds
\end{important}

\begin{formula}{t-SNE (t-Distributed Stochastic Neighbor Embedding)}
\textbf{Purpose:} Non-linear dimensionality reduction for visualization

\textbf{High-dimensional similarity:}
\[p_{j|i} = \frac{\exp(-\|x_{i} - x_j\|^2 / 2\sigma_i^2)}{\sum_{k \neq i}\exp(-\|x_{i} - x_k\|^2 / 2\sigma_i^2)}\]

\textbf{Low-dimensional similarity (t-distribution):}
\[q_{ij} = \frac{(1 + \|y_{i} - y_j\|^2)^{-1}}{\sum_{k \neq l}(1 + \|y_k - y_l\|^2)^{-1}}\]

\textbf{Cost Function (KL divergence):}
\[C = \sum_i KL(P_i \| Q_i) = \sum_i \sum_j p_{ij}\log\frac{p_{ij}}{q_{ij}}\]

\textbf{When to use:} Visualizing high-dimensional data (2D/3D plots)

\textbf{Note:} Primarily for visualization, not for dimensionality reduction in modeling
\end{formula}

% ============================================================
% MARKOV CHAINS
% ============================================================

\section{Markov Chains}

\begin{important}{Markov Chain Basics - MEDIUM PROBABILITY}
\textbf{Definition:} Stochastic process where future state depends only on current state (memoryless).

\textbf{Markov Property:}
\[P(X_{n+1} = j | X_n = i, X_{n-1}, \ldots, X_0) = P(X_{n+1} = j | X_n = i)\]

\textbf{Transition Probability:}
\[P_{ij} = P(X_{n+1} = j | X_n = i)\]
\end{important}

\begin{important}{Transition Matrix - HIGH PROBABILITY}
\textbf{Matrix P:}
\[\mathbf{P} = \begin{pmatrix}
P_{11} & P_{12} & \cdots & P_{1n} \\
P_{21} & P_{22} & \cdots & P_{2n} \\
\vdots & \vdots & \ddots & \vdots \\
P_{n1} & P_{n2} & \cdots & P_{nn}
\end{pmatrix}\]

\textbf{Properties:}
- Row sums = 1: \(\sum_{j=1}^{n} P_{ij} = 1\)
- All entries $\geq$ 0

\textbf{n-step Transition Probability:}
\[P^{(n)}_{ij} = (\mathbf{P}^n)_{ij}\]

\textbf{State Distribution after n steps:}
\[\boldsymbol{\pi}_n = \boldsymbol{\pi}_0 \mathbf{P}^n\]
\end{important}

\begin{important}{Stationary Distribution - MEDIUM PROBABILITY}
\textbf{Definition:} Distribution $\pi$ that doesn't change under transition
\[\boldsymbol{\pi} = \boldsymbol{\pi}\mathbf{P}\]

\textbf{Or equivalently:}
\[\boldsymbol{\pi}^T\mathbf{P} = \boldsymbol{\pi}^T\]

\textbf{With constraint:}
\[\sum_{i=1}^{n} \pi_i = 1\]

\textbf{Interpretation:} Long-run equilibrium probabilities.

\textbf{Note:} For ergodic chains, stationary distribution is unique and \(\lim_{n \to \infty} \mathbf{P}^n\) converges to $\pi$.

\textbf{Computing Stationary Distribution:}
Solve: (P^T - I)$\pi$ = 0 subject to $\Sigma$$\pi$_i = 1
\end{important}

\begin{formula}{Markov Chain Properties}
\textbf{Irreducibility:} Can reach any state from any other state

\textbf{Aperiodicity:} gcd of return times = 1

\textbf{Ergodicity:} Irreducible + aperiodic
- Guarantees unique stationary distribution
- Time-average = ensemble average

\textbf{Detailed Balance:}
\[\pi_i P_{ij} = \pi_j P_{ji}\]

If satisfied, $\pi$ is stationary distribution (reversible chain)

\textbf{Mean Recurrence Time:}
\[\mu_i = \frac{1}{\pi_i}\]

Average time to return to state i
\end{formula}

\begin{formula}{PageRank Algorithm}
\textbf{Application of Markov Chains:} Google's PageRank

\textbf{Formula:}
\[PR(A) = (1-d) + d\sum_{i=1}^{n}\frac{PR(T_{i})}{C(T_{i})}\]

where:
- d = damping factor (typically 0.85)
- T_{i} = pages linking to A
- C(T_{i}) = number of outbound links from T_{i}

\textbf{Matrix Form:}
\[\mathbf{PR} = (1-d)\mathbf{e} + d\mathbf{M}\mathbf{PR}\]

\textbf{Interpretation:} Stationary distribution of random web surfer
\end{formula}

% ============================================================
% MONTE CARLO METHODS
% ============================================================

\section{Monte Carlo Methods}

\begin{formula}{Monte Carlo Integration - LOW PROBABILITY}
\textbf{Estimate integral:}
\[\int_a^b f(x)\,dx \approx \frac{b-a}{n}\sum_{i=1}^{n} f(x_{i})\]

where x_{i} are random samples from uniform distribution on [a, b].

\textbf{Expected Value:}
\[\mathbb{E}[g(X)] \approx \frac{1}{n}\sum_{i=1}^{n} g(x_{i})\]

where x_{i} are samples from distribution of X.
\end{formula}

\begin{formula}{Monte Carlo Simulation - MEDIUM PROBABILITY}
\textbf{Steps:}
1. Define probability distributions for inputs
2. Generate random samples from distributions
3. Run deterministic computation for each sample
4. Aggregate results (mean, variance, percentiles)

\textbf{Law of Large Numbers:} As n → $\infty$, sample mean converges to true mean.

\textbf{Standard Error:}
\[SE = \frac{s}{\sqrt{n}}\]

where s is sample standard deviation.

\textbf{Confidence Interval:}
\[\bar{x} \pm 1.96 \cdot \frac{s}{\sqrt{n}}\]

\textbf{Sample Size for Desired Accuracy:}
\[n = \left(\frac{z_{\alpha/2} \cdot s}{\epsilon}\right)^2\]

where $\varepsilon$ is desired margin of error.
\end{formula}

\begin{formula}{Importance Sampling}
\textbf{For rare events:} Sample from importance distribution

\textbf{Estimator:}
\[\hat{\theta} = \frac{1}{n}\sum_{i=1}^{n}\frac{f(x_{i})}{g(x_{i})}h(x_{i})\]

where:
- f(x) = target distribution
- g(x) = importance distribution (easier to sample)
- h(x) = function of interest

\textbf{Variance Reduction:} Can significantly reduce MC variance

\textbf{When to use:} Rare events, tail probabilities, difficult distributions
\end{formula}

\begin{formula}{Markov Chain Monte Carlo (MCMC)}
\textbf{Purpose:} Sample from complex distributions

\textbf{Metropolis-Hastings Algorithm:}
1. Start with x$_0$
2. Propose x* from q(x*|x_{t-1})
3. Compute acceptance probability:
\[\alpha = \min\left(1, \frac{p(x^*)q(x_{t-1}|x^*)}{p(x_{t-1})q(x^*|x_{t-1})}\right)\]
4. Accept x* with probability $\alpha$, otherwise keep x_{t-1}
5. Repeat

\textbf{Gibbs Sampling:} Special case where proposals always accepted

\textbf{When to use:} Bayesian inference, complex posteriors
\end{formula}

% ============================================================
% DETERMINISTIC MODELS
% ============================================================

\section{Deterministic Models}

\begin{formula}{Deterministic vs Stochastic}
\textbf{Deterministic Model:}
- No randomness
- Same input always gives same output
- Described by equations

\textbf{Examples:}
- Linear equations: y = mx + b
- Differential equations: dy/dt = f(t, y)
- Physical laws: F = ma

\textbf{Stochastic Model:}
- Includes randomness
- Same input can give different outputs
- Uses probability distributions

\textbf{Examples:}
- Monte Carlo simulations
- Markov chains
- Random walk
\end{formula}

\begin{formula}{Common Deterministic Models}
\textbf{Linear Growth:}
\[y(t) = y_{0} + rt\]

\textbf{Exponential Growth:}
\[y(t) = y_{0} e^{rt}\]

\textbf{Logistic Growth:}
\[y(t) = \frac{K}{1 + \left(\frac{K-y_{0}}{y_{0}}\right)e^{-rt}}\]

where:
- y_{0} = initial value
- r = growth rate
- K = carrying capacity (for logistic)
\end{formula}

\subsection{Common Deterministic Models}
\textbf{Linear Growth:}
\[y(t) = y_{0} + rt\]

\textbf{Exponential Growth:}
\[y(t) = y_{0} e^{rt}\]

\textbf{Logistic Growth:}
\[y(t) = \frac{K}{1 + \left(\frac{K-y_{0}}{y_{0}}\right)e^{-rt}}\]

where:
- y_{0} = initial value
- r = growth rate
- K = carrying capacity (for logistic)
\end{formula}

% ============================================================
% BIAS-VARIANCE TRADEOFF
% ============================================================

\section{Bias-Variance Tradeoff}

\begin{important}{Bias-Variance Decomposition - MEDIUM PROBABILITY}
\textbf{Expected Prediction Error:}
\[\mathbb{E}[(y - \hat{f}(x))^2] = \text{Bias}^2 + \text{Variance} + \text{Irreducible Error}\]

\textbf{Bias:}
\[\text{Bias}[\hat{f}(x)] = \mathbb{E}[\hat{f}(x)] - f(x)\]

Measures how far off average predictions are from true value.

\textbf{Variance:}
\[\text{Var}[\hat{f}(x)] = \mathbb{E}[(\hat{f}(x) - \mathbb{E}[\hat{f}(x)])^2]\]

Measures variability of predictions across different training sets.

\textbf{Tradeoff:}
- High bias, low variance: Underfitting (too simple)
- Low bias, high variance: Overfitting (too complex)
- Goal: Find balance that minimizes total error
\end{important}

\begin{example}
\textbf{Model Complexity:}
- Simple models: High bias, low variance (linear regression)
- Complex models: Low bias, high variance (high-degree polynomial)
- Ensemble methods: Balance both (random forest, boosting)

\textbf{Regularization:} Increases bias, decreases variance
\end{example}

% ============================================================
% FEATURE SELECTION
% ============================================================

\section{Feature Selection and Engineering}

\begin{formula}{Filter Methods}
\textbf{Correlation with Target:}
Select features with highest |correlation| with target variable.

\textbf{Chi-Square Test:}
For categorical features, test independence with target.

\textbf{Mutual Information:}
\[I(X;Y) = \sum_{x,y} p(x,y)\log\frac{p(x,y)}{p(x)p(y)}\]

Measures dependency between features and target.

\textbf{Variance Threshold:}
Remove features with low variance (nearly constant).
\end{formula}

\begin{formula}{Wrapper Methods}
\textbf{Forward Selection:}
1. Start with no features
2. Add feature that improves model most
3. Repeat until no improvement

\textbf{Backward Elimination:}
1. Start with all features
2. Remove feature that hurts least
3. Repeat until further removal hurts

\textbf{Recursive Feature Elimination (RFE):}
Train model, rank features by importance, remove least important, repeat.

\textbf{When to use:} Better performance than filters, but computationally expensive.
\end{formula}

\begin{formula}{Embedded Methods}
\textbf{LASSO:} L1 penalty automatically selects features (sets coefficients to zero)

\textbf{Ridge:} L2 penalty shrinks but doesn't eliminate features

\textbf{Elastic Net:} Combination of L1 and L2

\textbf{Tree-Based Feature Importance:}
Based on information gain or impurity decrease.

\textbf{Advantage:} Feature selection during model training.
\end{formula}

% ============================================================
% ENSEMBLE METHODS
% ============================================================

\section{Ensemble Methods}

\begin{important}{Bagging (Bootstrap Aggregating) - MEDIUM PROBABILITY}
\textbf{Algorithm:}
1. Create B bootstrap samples (sampling with replacement)
2. Train model on each sample
3. Aggregate predictions (average or vote)

\textbf{Purpose:} Reduce variance, prevent overfitting

\textbf{Examples:} Random Forest, Bagged Decision Trees

\textbf{Prediction:}
\[\hat{y} = \frac{1}{B}\sum_{b=1}^{B}\hat{f}_b(x)\]

\textbf{When to use:} High-variance models (decision trees, neural networks)
\end{important}

\begin{important}{Boosting - MEDIUM PROBABILITY}
\textbf{Idea:} Sequentially train models, each correcting previous errors

\textbf{AdaBoost Weight Update:}
\[w_i^{(t+1)} = w_i^{(t)} \exp(\alpha_t \mathbb{I}(y_{i} \neq \hat{y}_i))\]

where:
\[\alpha_t = \frac{1}{2}\log\frac{1 - \\varepsilon_{t}}{\\varepsilon_{t}}\]
$\varepsilon$_t = weighted error rate

\textbf{Gradient Boosting:}
Fit new model to residuals (negative gradient) of previous model.

\textbf{Purpose:} Reduce bias, improve accuracy

\textbf{Examples:} AdaBoost, XGBoost, LightGBM, CatBoost

\textbf{When to use:} Need high accuracy, can handle complex patterns
\end{important}

\begin{formula}{Stacking (Stacked Generalization)}
\textbf{Two-Level Approach:}

\textbf{Level 0 (Base Models):}
Train multiple diverse models on training data.

\textbf{Level 1 (Meta-Model):}
Train on predictions from base models.

\textbf{Prediction:}
\[\hat{y} = g(f_1(x), f_2(x), \ldots, f_M(x))\]

where g is meta-model, f_i are base models.

\textbf{When to use:} Combining different model types for best performance.
\end{formula}

% ============================================================
% TIME SERIES BASICS
% ============================================================

\section{Time Series Analysis (Brief)}

\begin{formula}{Moving Average (MA)}
\textbf{Simple Moving Average:}
\[MA_t = \frac{1}{k}\sum_{i=0}^{k-1}y_{t-i}\]

\textbf{Weighted Moving Average:}
\[WMA_t = \sum_{i=0}^{k-1}w_i y_{t-i}\]

\textbf{Exponential Smoothing:}
\[S_t = \alpha y_{t} + (1-\alpha)S_{t-1}\]

where 0 < $\alpha$ < 1 is smoothing parameter.

\textbf{When to use:} Smoothing data, forecasting, trend detection.
\end{formula}

\begin{formula}{Autocorrelation}
\textbf{Autocorrelation Function (ACF):}
\[\rho_k = \frac{\text{Cov}(y_{t}, y_{t-k})}{\text{Var}(y_{t})}\]

Correlation between series and itself at lag k.

\textbf{Partial Autocorrelation (PACF):}
Correlation between y_{t} and y_{t-k} after removing effect of intermediate lags.

\textbf{When to use:} Identifying AR/MA orders, checking randomness.
\end{formula}

% ============================================================
% QUICK REFERENCE TABLES
% ============================================================

\section{Quick Reference Tables}

\subsection{Distribution Properties}

\begin{center}
\begin{tabular}{|l|c|c|c|}
\hline
\textbf{Distribution} & \textbf{E[X]} & \textbf{Var(X)} & \textbf{Use Case} \\
\hline
Bernoulli(p) & p & p(1-p) & Single trial \\
Binomial(n,p) & np & np(1-p) & n trials \\
Poisson($\lambda$) & $\lambda$ & $\lambda$ & Rare events \\
Normal($\mu$,$\sigma$²) & $\mu$ & $\sigma$² & Continuous, symmetric \\
\hline
\end{tabular}
\end{center}

\subsection{Test Selection Guide}

\begin{center}
\begin{tabular}{|p{5cm}|p{8cm}|}
\hline
\textbf{Scenario} & \textbf{Test} \\
\hline
One sample, known $\sigma$ & Z-test \\
One sample, unknown $\sigma$ & t-test \\
Two samples, independent & Two-sample t-test \\
Two samples, paired & Paired t-test \\
Proportions & Z-test for proportions \\
\hline
\end{tabular}
\end{center}

\subsection{Common Z-scores}

\begin{center}
\begin{tabular}{|c|c|c|}
\hline
\textbf{Confidence Level} & \textbf{$\alpha$} & \textbf{Z-score} \\
\hline
90\% & 0.10 & 1.645 \\
95\% & 0.05 & 1.96 \\
99\% & 0.01 & 2.576 \\
99.9\% & 0.001 & 3.291 \\
\hline
\end{tabular}
\end{center}

\subsection{Effect Size Interpretations}

\begin{center}
\small
\begin{tabular}{|l|c|c|c|}
\hline
\textbf{Measure} & \textbf{Small} & \textbf{Medium} & \textbf{Large} \\
\hline
Cohen's d & 0.2 & 0.5 & 0.8 \\
Pearson r & 0.1 & 0.3 & 0.5 \\
R² & 0.01 & 0.09 & 0.25 \\
Cramér's V & 0.1 & 0.3 & 0.5 \\
Eta-squared ($\eta$²) & 0.01 & 0.06 & 0.14 \\
\hline
\end{tabular}
\end{center}

\subsection{Distance Metrics Comparison}

\begin{center}
\small
\begin{tabular}{|l|p{6cm}|p{5cm}|}
\hline
\textbf{Distance} & \textbf{Formula} & \textbf{Use Case} \\
\hline
Euclidean & $\sqrt{\sum(x_{i} - y_{i})^2}$ & General purpose, KNN \\
Manhattan & $\sum|x_{i} - y_{i}|$ & Grid-like paths, robust \\
Cosine & $1 - \frac{x \cdot y}{\|x\|\|y\|}$ & Text, high dimensions \\
Minkowski (p=3) & $(\sum|x_{i} - y_{i}|^p)^{1/p}$ & Generalization \\
Hamming & Count of differences & Binary/categorical \\
\hline
\end{tabular}
\end{center}

\subsection{Model Complexity and Regularization}

\begin{center}
\small
\begin{tabular}{|l|c|c|c|}
\hline
\textbf{Method} & \textbf{Bias} & \textbf{Variance} & \textbf{Sparsity} \\
\hline
OLS & Low & High & No \\
Ridge (L2) & Medium & Medium & No \\
LASSO (L1) & Medium & Medium & Yes \\
Elastic Net & Medium & Medium & Yes \\
\hline
\end{tabular}
\end{center}

\subsection{Evaluation Metrics by Task}

\begin{center}
\small
\begin{tabular}{|l|p{10cm}|}
\hline
\textbf{Task} & \textbf{Recommended Metrics} \\
\hline
Regression & RMSE, MAE, R², Adjusted R² \\
Binary Classification & Accuracy, Precision, Recall, F1, AUC-ROC \\
Multi-class Classification & Accuracy, Macro/Micro F1, Confusion Matrix \\
Imbalanced Classification & F1-Score, Precision-Recall AUC, MCC \\
Clustering & Silhouette Score, Davies-Bouldin Index, WCSS \\
Ranking & NDCG, MAP, MRR \\
\hline
\end{tabular}
\end{center}

\subsection{Algorithm Comparison}

\begin{center}
\small
\begin{tabular}{|l|c|c|c|}
\hline
\textbf{Algorithm} & \textbf{Type} & \textbf{Pros} & \textbf{Cons} \\
\hline
KNN & Supervised & Simple, no training & Slow prediction, memory \\
K-Means & Unsupervised & Fast, scalable & Need k, spherical clusters \\
Naive Bayes & Supervised & Fast, works with small data & Independence assumption \\
Decision Tree & Supervised & Interpretable & Overfitting \\
\hline
\end{tabular}
\end{center}

% ============================================================
% ADDITIONAL DISTRIBUTIONS
% ============================================================

\section{Additional Important Distributions}

\begin{important}{T-Distribution - VERY HIGH PROBABILITY}
\textbf{Probability Density Function:}
\[f(t) = \frac{\Gamma\left(\frac{\nu+1}{2}\right)}{\sqrt{\nu\pi}\,\Gamma\left(\frac{\nu}{2}\right)}\left(1 + \frac{t^2}{\nu}\right)^{-(\nu+1)/2}, \quad t \in \mathbb{R}\]

where $\nu$ = degrees of freedom, $\Gamma(\cdot)$ = Gamma function

\textbf{Properties:}
\begin{itemize}[nosep]
\item $\mathbb{E}[T] = 0$ for $\nu > 1$
\item $\text{Var}(T) = \frac{\nu}{\nu-2}$ for $\nu > 2$
\item As $\nu \to \infty$, $t_\nu \to \mathcal{N}(0,1)$
\item Heavier tails than normal (more robust to outliers)
\item Symmetric about 0
\end{itemize}

\textbf{Common Critical Values:}
$t_{0.025,10} = 2.228$, $t_{0.025,20} = 2.086$, $t_{0.025,30} = 2.042$, $t_{0.025,\infty} = 1.96$
\end{important}

\begin{important}{Chi-Square Distribution $\chi^2(\nu)$ - HIGH PROBABILITY}
\textbf{Probability Density Function:}
\[f(x) = \frac{1}{2^{\nu/2}\Gamma(\nu/2)}x^{(\nu/2)-1}e^{-x/2}, \quad x > 0\]

\textbf{Properties:}
$\mathbb{E}[\chi^2] = \nu$, $\text{Var}(\chi^2) = 2\nu$

If $Z_1, \ldots, Z_\nu \overset{\text{iid}}{\sim} \mathcal{N}(0,1)$, then $\sum_{i=1}^\nu Z_i^2 \sim \chi^2(\nu)$

\textbf{Variance Estimation:} $\frac{(n-1)s^2}{\sigma^2} \sim \chi^2(n-1)$
\end{important}

\begin{formula}{F-Distribution $F(\nu_1, \nu_2)$ - MEDIUM PROBABILITY}
If $U \sim \chi^2(\nu_1)$ and $V \sim \chi^2(\nu_2)$ independent:
\[F = \frac{U/\nu_1}{V/\nu_2} \sim F(\nu_1, \nu_2)\]

\textbf{Properties:} $\mathbb{E}[F] = \frac{\nu_2}{\nu_2-2}$ for $\nu_2 > 2$

\textbf{Comparing Variances:} $\frac{s_1^2}{s_2^2} \sim F(n_1-1, n_2-1)$
\end{formula}

\begin{important}{Gamma Distribution $\Gamma(\alpha, \beta)$ - MEDIUM PROBABILITY}
\[f(x) = \frac{\beta^\alpha}{\Gamma(\alpha)}x^{\alpha-1}e^{-\beta x}, \quad x > 0\]

\textbf{Parameters:} $\alpha$ = shape, $\beta$ = rate

\textbf{Properties:} $\mathbb{E}[X] = \frac{\alpha}{\beta}$, $\text{Var}(X) = \frac{\alpha}{\beta^2}$

\textbf{Special Cases:} $\Gamma(1, \beta) = \text{Exp}(\beta)$, $\Gamma(\nu/2, 1/2) = \chi^2(\nu)$
\end{important}

% ============================================================
% SAMPLING DISTRIBUTIONS
% ============================================================

\section{Sampling Distributions and Convergence}

\begin{important}{Sampling Distribution of $\bar{X}$ - VERY HIGH PROBABILITY}
For i.i.d. $X_1, \ldots, X_n$ with $\mathbb{E}[X_i] = \mu$, $\text{Var}(X_i) = \sigma^2$:

\textbf{Mean and Variance:}
\[\mathbb{E}[\bar{X}] = \mu, \quad \text{Var}(\bar{X}) = \frac{\sigma^2}{n}, \quad \text{SE}(\bar{X}) = \frac{\sigma}{\sqrt{n}}\]

\textbf{Distribution:}
\[\bar{X} \sim \mathcal{N}\left(\mu, \frac{\sigma^2}{n}\right) \quad \text{(if $X_i$ normal or $n$ large by CLT)}\]

\textbf{Standardized:}
\[Z = \frac{\bar{X} - \mu}{\sigma/\sqrt{n}} \sim \mathcal{N}(0,1) \quad \text{or} \quad T = \frac{\bar{X} - \mu}{s/\sqrt{n}} \sim t_{n-1}\]
\end{important}

\begin{formula}{Chi-Square for Variance - HIGH PROBABILITY}
For $X_1, \ldots, X_n \sim \mathcal{N}(\mu, \sigma^2)$:
\[\frac{(n-1)s^2}{\sigma^2} \sim \chi^2(n-1)\]

\textbf{CI for $\sigma^2$:}
\[\left[\frac{(n-1)s^2}{\chi^2_{\alpha/2,n-1}}, \frac{(n-1)s^2}{\chi^2_{1-\alpha/2,n-1}}\right]\]
\end{formula}

% ============================================================
% NON-PARAMETRIC METHODS
% ============================================================

\section{Non-Parametric Methods}

\begin{formula}{Kernel Density Estimation (KDE)}
\[\hat{f}(x) = \frac{1}{nh}\sum_{i=1}^{n}K\left(\frac{x - x_i}{h}\right)\]

where $K(\cdot)$ = kernel, $h$ = bandwidth

\textbf{Gaussian Kernel:} $K(u) = \frac{1}{\sqrt{2\pi}}e^{-u^2/2}$

\textbf{Silverman's Bandwidth:} $h = 0.9 \cdot \min\left\{s, \frac{\text{IQR}}{1.34}\right\} \cdot n^{-1/5}$
\end{formula}

\begin{important}{Bootstrap Resampling - MEDIUM PROBABILITY}
1. Draw $B$ samples of size $n$ with replacement
2. Calculate $\hat{\theta}^{(b)}$ for each sample
3. Estimate SE and CI

\textbf{Bootstrap SE:}
\[\text{SE}_{\text{boot}} = \sqrt{\frac{1}{B-1}\sum_{b=1}^{B}\left(\hat{\theta}^{(b)} - \bar{\theta}^*\right)^2}\]

\textbf{Percentile CI:} $[\hat{\theta}_{(0.025)}, \hat{\theta}_{(0.975)}]$
\end{important}

% ============================================================
% CURSE OF DIMENSIONALITY
% ============================================================

\section{High-Dimensional Statistics}

\begin{important}{Curse of Dimensionality - MEDIUM PROBABILITY}
\textbf{Volume of Hypersphere:}
\[V_d(r) = \frac{\pi^{d/2}}{\Gamma(d/2 + 1)}r^d\]

\textbf{Key Issues:}
\begin{itemize}[nosep]
\item Distance concentration: $\frac{d_{\max} - d_{\min}}{d_{\min}} \to 0$ as $d \to \infty$
\item Most volume near surface
\item Data becomes sparse
\item $n$ required grows exponentially with $d$
\end{itemize}

\textbf{Mitigation:}
- Dimensionality reduction (PCA)
- Feature selection
- Regularization
- Distance metric learning
\end{important}

% ============================================================
% INFORMATION THEORY BASICS
% ============================================================

\section{Information Theory for ML}

\begin{important}{Entropy - HIGH PROBABILITY}
\textbf{Shannon Entropy (Binary):}
\[H(p) = -p\log_2 p - (1-p)\log_2(1-p)\]

\textbf{General Entropy:}
\[H(X) = -\sum_{i=1}^{k}p_i\log_2 p_i = \mathbb{E}[-\log_2 P(X)]\]

\textbf{Properties:}
\begin{itemize}[nosep]
\item $H(X) \geq 0$ (non-negative)
\item $H(X) = 0$ iff distribution is deterministic
\item Maximum when uniform: $H_{\max} = \log_2 k$
\item Concave function
\end{itemize}
\end{important}

\begin{important}{Gini Impurity - HIGH PROBABILITY}
\[G(p) = 2p(1-p) = 1 - \sum_{i=1}^{k}p_i^2\]

\textbf{Relationship:} $H(p) \geq G(p)$ for all $p$

\textbf{For Decision Trees:}
- Information Gain: $\text{IG} = H(S) - \sum_{v}\frac{|S_v|}{|S|}H(S_v)$
- Gini Gain: $\Delta G = G(S) - \sum_{v}\frac{|S_v|}{|S|}G(S_v)$
\end{important}

\begin{formula}{Cross-Entropy Loss}
\textbf{Binary Classification:}
\[L = -\frac{1}{n}\sum_{i=1}^{n}\left[y_i\log(\hat{p}_i) + (1-y_i)\log(1-\hat{p}_i)\right]\]

\textbf{Multi-Class:}
\[L = -\frac{1}{n}\sum_{i=1}^{n}\sum_{c=1}^{C}y_{ic}\log(\hat{p}_{ic})\]

\textbf{Interpretation:} Measure of difference between two probability distributions
\end{formula}

% ============================================================
% ADVANCED REGRESSION TOPICS
% ============================================================

\section{Advanced Regression Concepts}

\begin{important}{Residual Analysis - HIGH PROBABILITY}
\textbf{Residual:} $e_i = y_i - \hat{y}_i$

\textbf{Standardized Residual:}
\[r_i = \frac{e_i}{\hat{\sigma}\sqrt{1-h_{ii}}}\]

where $h_{ii} = [\mathbf{H}]_{ii}$ (leverage)

\textbf{Studentized Residual:}
\[t_i = \frac{e_i}{\hat{\sigma}_{(i)}\sqrt{1-h_{ii}}} \sim t_{n-p-2}\]

\textbf{Cook's Distance (Influence):}
\[D_i = \frac{(\hat{\boldsymbol{\beta}} - \hat{\boldsymbol{\beta}}_{(i)})^T(\mathbf{X}^T\mathbf{X})(\hat{\boldsymbol{\beta}} - \hat{\boldsymbol{\beta}}_{(i)})}{p\hat{\sigma}^2} = \frac{r_i^2}{p}\frac{h_{ii}}{1-h_{ii}}\]

\textbf{Rule:} $D_i > \frac{4}{n}$ or $D_i > 1$ suggests influential point
\end{important}

\begin{formula}{Variance Inflation Factor (VIF)}
For predictor $X_j$:
\[VIF_j = \frac{1}{1-R_j^2}\]

where $R_j^2$ = R-squared from regressing $X_j$ on other predictors

\textbf{Interpretation:}
- $VIF < 5$: No multicollinearity
- $5 \leq VIF < 10$: Moderate
- $VIF \geq 10$: Severe multicollinearity
\end{formula}

\begin{formula}{Akaike Information Criterion (AIC)}
\[AIC = 2k - 2\ln(\hat{L})\]

where $k$ = number of parameters, $\hat{L}$ = maximized likelihood

\textbf{For Linear Regression:}
\[AIC = n\ln\left(\frac{SSE}{n}\right) + 2k\]

\textbf{Bayesian IC (BIC):}
\[BIC = k\ln(n) - 2\ln(\hat{L})\]

\textbf{Rule:} Lower AIC/BIC indicates better model (balances fit vs complexity)
\end{formula}

% ============================================================
% RESAMPLING METHODS
% ============================================================

\section{Cross-Validation and Model Selection}

\begin{important}{K-Fold Cross-Validation - HIGH PROBABILITY}
\textbf{Procedure:}
1. Split data into $K$ folds
2. For each fold $k=1,\ldots,K$:
   - Train on $K-1$ folds
   - Validate on fold $k$
   - Compute error $E_k$
3. Average: $CV_{(K)} = \frac{1}{K}\sum_{k=1}^{K}E_k$

\textbf{Special Cases:}
- $K=n$: Leave-One-Out CV (LOOCV)
- $K=5$ or $K=10$: Common choices

\textbf{Bias-Variance:}
- Small $K$ → higher bias, lower variance
- Large $K$ → lower bias, higher variance
\end{important}

\begin{formula}{Train-Test Split}
\textbf{Common Splits:}
- 70/30 train/test
- 80/20 train/test
- 60/20/20 train/validation/test

\textbf{Stratified Sampling:} Preserve class proportions in splits
\end{formula}

\vspace{1cm}
\begin{center}
\textbf{\large --- END OF FORMULA SHEET ---}

\textit{Good luck on your exam!}
\end{center}

\end{document}
